\section{Scope, prior work, and the proved contributions}
\label{sec:scope}

The source manuscript is \emph{Polylogarithms and their Arithmetic
Bridges}, in the directory
\src{Analysis/Polylogarithms/docs/manuscript} of ProveIt
\cite{repoManuscript}. We fixed the inspected revision at
\begin{center}
\small\src{d00eba30c1ce73ee5cfbeb50bde9095246b9c25c}.
\end{center}
That revision contains the consolidated manuscript, its earlier reports,
and four relevant incoming archives. The present article is an additional
research contribution prepared for integration; it is not a replacement
for the whole manuscript or a certification of every historical claim.

The immediately preceding reports have already proved Gaussian double
parity formulas, the weight-five shuffle relation, several one-index-two
multiple-polylogarithm reductions, and the cubic and quartic
complementary-power trilogarithm examples. They also reduce the four
same-point weight-five Gaussian coordinates to at most two modulo
specified single-value products. We use those facts where appropriate
and do not count their rederivation as a new contribution. In particular,
the first follow-up problem in the exact-reductions package is the
specific $S_4$ relation proved below; its fifth asks for a classification
of complementary-power redundancies \cite{repoExact}.

\subsection{A precise contribution ledger}

\begin{longtable}{@{}>{\raggedright\arraybackslash}p{0.28\textwidth}>{\raggedright\arraybackslash}p{0.66\textwidth}@{}}
\toprule
Result & Mathematical content and proof status\\
\midrule\endhead
Complementary-power rigidity & A complete theorem for two positive
integer exponent pairs at the same real base, obtained from Ljunggren's
trinomial theorem and a root-rotation argument. It includes equal
exponents and arbitrary common divisors.\\[3pt]
Signed cyclic seeds & Exact cardinalities $0,3,6,12$, with the case
twelve characterized; a bound $b\leq\lfloor5D/3\rfloor$ for bases of
degree at most $D$. These are statements about pure powers.\\[3pt]
Plastic ladders & Analytic proofs of two specific manuscript identities,
plus explicit dilogarithm and trilogarithm identities with no
logarithmic terms. Classical functional equations are the inputs.\\[3pt]
Gaussian $S_4$ & A proved identity with a finite exact certificate.
The new ingredient relative to the preceding attempts is an explicit
M\"obius duality relation.\\[3pt]
Proportional harmonic depth & An explicit positive, decreasing, convex
exponential rate and a computable first correction, with an all-orders
expansion uniform on compact parameter sets.\\[3pt]
Reflected gamma moments & Exact poles and residue polynomials, shifted
factorial growth, late-coefficient asymptotics, and an explicit
truncation bound.\\[3pt]
Residue and log-sine identities & An exact finite reflection identity
and an unconditional formal polynomial version of a coefficient
conjecture recorded in a 2010 author preprint.\\
\bottomrule
\end{longtable}

The classification gives a particularly useful change in how one can
search. A box of polynomial coefficients is not exhaustive, even at
fixed degree. In contrast, the exponent bound proved here supplies an
exhaustive finite search for the restricted equation $r^a+r^b=1$.
Equally importantly, it identifies exactly what this search cannot
settle: multiplicative relations involving products of many factors
$1-r^j$, additional units, or more general polylogarithm arguments are
outside its scope.

The $S_4$ proof has a different character. It is a finite calculation
with convergent iterated-integral values, using explicitly justified
single-divergence relations as well as convergent identities. Its exact certificate can be checked by
rational arithmetic. The mathematical conclusion is an equality of
the corresponding periods. It is not a proof that a computed quotient
of all supplied relations is the full space of period relations, nor
does it prove that any remaining coordinate is irreducible.

\subsection{Classical inputs and attribution}

Ljunggren's irreducibility theorem is used in the precise form recorded
by Boyd, Martin, and Thom \cite[Lemma 4.1]{BoydMartinThom2015}; Tverberg's
proof is another primary reference \cite{Tverberg1960}. Rogers' pentagon
and Kummer's trilogarithm equation are classical; the modified
trilogarithm normalization and its functional equations are taken from
Zagier \cite{ZagierPolylogs1991}. The shuffle and stuffle laws and their
regularization belong to the established theory of multiple
polylogarithms \cite{jqzhao2010,PanzerParity}. The general existence
of parity reductions is therefore not being claimed here as new.

The analytic calculations use standard Mellin continuation, Lagrange
inversion, and Laplace asymptotics. What is derived is their explicit
application to the stated families, including normalization constants,
finite coefficient identities, and error terms. Mixed reflected
log-gamma moments were already introduced by Amdeberhan and coauthors
\cite{ACEMKM}; neither their definition nor their reflection symmetry
is a new observation. Their large-exponent and late-coefficient
behavior is the subject of the later sections.

This article establishes the statements it proves without asserting
first discovery in the entire literature. Exact priority for equivalent
normalizations of special-value identities is a separate bibliographic
question. The useful distinction is between proved results, numerical
checks of those results, and questions that remain open.

\subsection{Conventions}

For positive integers $s_1,\ldots,s_d$, our multiple polylogarithm uses
decreasing indices:
\begin{equation}\label{eq:mpl-convention}
 \Li_{s_1,\ldots,s_d}(z_1,\ldots,z_d)
 =\sum_{n_1>\cdots>n_d\geq1}
 \frac{z_1^{n_1}\cdots z_d^{n_d}}
 {n_1^{s_1}\cdots n_d^{s_d}}.
\end{equation}
Boundary values are assigned by the ordinary convergent series or by
the explicitly stated regularization. The corresponding hyperlogarithm
uses $dt/(t-a)$ and an outer-first word, so its integration simplex is
$0<t_N<\cdots<t_1<1$. The word letters associated to
\eqref{eq:mpl-convention} are reciprocal \emph{prefix} products.
These choices are essential to the signs in the Gaussian certificate.

All logarithms in the real algebraic and gamma sections are real
logarithms of positive numbers. Complex arguments in the Gaussian
section are connected to their convergent series, fixing the branches.
We write $G=\beta(2)$ for Catalan's constant and
$\beta(s)=\sum_{j\geq0}(-1)^j/(2j+1)^s$. The letter $h$ denotes
harmonic depth in Section~\ref{sec:saddle}; the fixed reflected exponent
in the gamma section is $m$. Repeated use of an auxiliary symbol in a
later, explicitly defined family does not identify the corresponding
constants across sections.
